% Conclusion --- derived from the Avance7 conclusions, restated with the
% corrected final numbers and defensible framing.
\section{Conclusion}
\label{sec:conclusion}

AgroSatCopilot shows that crop-area quantification from Sentinel-2 time
series is well served by a perceiver/reasoner design. Specialised
perception models --- temporal transformers, an XGBoost classifier over
AlphaEarth annual embeddings~\cite{brown2025alphaearth}, and a
contrastive vision-language branch~\cite{li2025farslip} --- are combined
through heterogeneous stacking, which improves the out-of-fold macro F1
over any individual model. The best individual perceiver, TSViT-pheno,
reaches mIoU 0.6253 and macro F1 0.7500 on the held-out fold, and the
five-member stacking ensemble attains an out-of-fold macro F1 of 0.6477.

The \emph{Be My Eyes} pattern~\cite{huang2025bemyeyes} establishes a
clean separation between the models that \emph{see} and a frozen
reasoner that \emph{communicates}, guaranteeing auditability and
resistance to hallucination, while a dual cloud/on-premises backend
preserves data sovereignty. Our transfer study is deliberately honest:
zero-shot generalisation from France to Catalonia fails, and few-shot
fine-tuning recovers mIoU 0.2468, consistent with reported limits on
the spatial transferability of foundation-model
embeddings~\cite{harvesting2026alphaearth}.

The main outstanding challenge is the six minority classes with low
F1, which require additional data or targeted augmentation rather than
more models of the same kind. Future work includes the full FarSLIP
phenological study (US-072), the multi-region label-space crosswalk and
$k$-shot curve (US-073), and the quantitative conversational-layer
benchmark on the H100 (US-069).
